\scanpage{72}
% archive foliation: 71

Si d'autre part $\underline G$ est un schéma en groupes fini étale sur $S$, et si $X'$ est muni d'une structure de $\underline G_X$-torseur au dessus de $X^*$, alors le stabilisateur \inserted{dans} $\underline G_{\deleted{\unsure{X}}}$ de la section $t'$ est un sous-schéma en groupes \inserted{$\underline G_{t'}$} de $\underline G$, \deleted{\unsure{schéma en groupes}} \inserted{qui} opère \inserted{fid.} sur \deleted{\unsure{comme si}} $T_{X'/S,t'}$ et dont l'image dans \deleted{$T_{X'/S,t'}$} $G_{mS}\simeq\underline{\operatorname{Aut}}_{\mathcal O_S}T_{X'/S,t'}$ est $\mu_{nS}$, d'où un monomorphisme canonique
\[
(3)\qquad \mu_{nS}\hookrightarrow\underline G.
\]
\margin{\uline{NB} On suppose \unsure{de même} que $\underline G$ opère sur $X'$, de façon que $X'^*=X'|X^*$ devienne un $\underline G_{X^*}$-torseur sur $X^*$.}

Par ailleurs, désignant par
\[
f^o:X^*\longrightarrow S
\]
le morphisme induit par $f$, le $\underline G_{X^*}$-torseur $X'^*$ définit une section
\[
(4)\qquad \unsure{\xi|\zeta}_{X'/S}\in H^0(S,R^1f^o_*f^{o*}(\underline G)).
\]
Tout ceci posé, on a :

\uline{Proposition} $(X',t')$ au dessus de $(X,t)$, avec l'opération de $\underline G$ dessus, est défini à isom unique près par la connaissance de (3)(4) et de (2). \deleted{\unsure{(ce dernier}} De façon précise,
